\section{Discussion}\label{sec:discussion}

\subsection{Why the sequences behave as they do}

The transient ratio of Eq.~\ref{eq:transient} falls when two consecutive frames place a moving trap at positions whose fields interfere with a phase difference. Three things control it: the shift of the pattern between frames, which Proposition~\ref{prop:blend} turns into a closed form; the phase of the trap in the two frames, which the lock prescribes and the induction only initializes; and the uniformity of each frame, which the optimization supplies. Linear phase interpolation controls the second and not the third, and it has the lowest ratio of the three sequences because its frames are not equalized: the per-frame spread of its frames is of the order of \lpiSigmaFrameMed{} percent. The per-frame method controls the third and not the second, and in these experiments the induced phases turn out to move by at most \accKimFiveStepMax{} rad between frames, which is small enough for the blend to stay above $0.9$ on every seed at \parFrames{} frames. The proposed sequence controls both, and that is the reason it improves on linear phase interpolation by \diffTwoMean{} in the mean of the ratio. It does not improve on the induced sequence, because the second control is already adequate there.

The direction of the remaining difference needs an explanation. The induced sequence has the higher dip although its phase step is about \stepRatio{} times larger. An explanation consistent with the data is that the transient ratio depends on the shift of the pattern between frames, which differs between the two methods because they follow different time profiles: the baselines move at constant velocity, the proposed sequence moves with a peak velocity of $15/8$ times the mean. The ablation supports this reading: with a linear profile the sequence reaches \ablTLinearMean{}, above the induced sequence at \accKimFiveTransMean{}, and with the profile of the method it is below it. The comparison of the two methods on the dip is therefore partly a comparison of two time profiles and not only of two ways of computing the frames. We report the comparison as run, with the profile that each method was defined with, and the result with the linear profile as an ablation of the proposed sequence.

\subsection{What the lock buys and what it costs}

The lock buys a controlled phase step. Without it the median of the largest step is \ablTNoLockStep{} rad, the value of the induced sequences; with it, \ablTFullStep{} rad, a reduction by a factor of \lockStepRatio{} at equal iterations. It also narrows the spread of outcomes across seeds: the standard deviation of the ratio is \ablTFullSd{} with the lock and \ablTNoLockSd{} without. Whether a smaller step is worth having depends on the experiment. The measurement of Knottnerus et al.\ of atom survival against a prescribed phase slip per hologram says that slips around $\pi$ lose atoms and slips of several tenths of $\pi$ do not~\cite{lpi}; the induced sequence stays at about \fpeval{round(\accKimFiveStepMax/pi,2)}$\pi$ on the acceptance task, so on this evidence it has no atom-survival problem to fix. The margin could matter in a setting where the induced phase moves more: with a coarser trajectory, a larger array whose frames are computed in fewer iterations, or a device whose phase response adds a slip. The experiments here do not reach such a setting, and the paper makes no claim about atoms.

The lock costs in two places. It makes the sequence depend on a path that must be consistent: when the end phases prescribe a path along which the hologram cannot equalize the traps, the optimization cannot repair it, and the seeds that fail on the long moves are of this kind (Conjecture~\ref{conj:phasepath}). The per-frame method has no such failure in these experiments, since it always follows the phase that the current frame prefers. And it costs iterations: the sequence at \parFrameIters{} iterations per frame needs \secRatio{} times the seconds of the induced sequence at its own count, although the ablation shows that 3 iterations suffice on the acceptance task.

\subsection{Failure cases and negative results}

The negative results are the following. The sequence does not exceed the induced sequence on the dip at 13 frames, nor at 7, and ties it at 25 and at equal work (Tables~\ref{tab:main} and~\ref{tab:frames}). With the frames of the bound it fails on \fpeval{\longSeeds-\longAMet} of \longSeeds{} seeds of each lattice move, by a wide margin, and the doubling of the frames does not rescue them, while the induced sequence reaches $0.9$ on all of them (Table~\ref{tab:long}). On the sort at 25 frames every method is below $0.9$, and the sequence is the lowest of the optimized ones (Fig.~\ref{fig:long}d). The lock removal does not bring the sequence below $0.9$ on a majority of seeds, so the lock is not shown to be what keeps the acceptance task above $0.9$. And the bound of the frame-count equation is not necessary for a short move, because the acceptance task is above $0.9$ at fewer frames than the bound asks. The value of the bound is a calibration on a single case, and a user who needs a guarantee for every seed cannot take it from this paper.

\subsection{Relation to the published methods}

Per-frame Gerchberg--Saxton with phase induction is the published method that computes every frame and keeps the phases continuous by initialization~\cite{kimoe}; the measurements here show that on a moving-trap task it is a strong baseline. The claim that survives the comparison is the narrow one stated in the introduction: the phase trajectory is prescribed and locked, and the number of frames has a derived bound, with the consequences measured above, a phase step that is smaller by a factor of \stepRatio{}, a per-frame spread that is smaller by a factor of \sigmaRatio{}, and a bound that is a calibration and fails on some seeds. The claim that the frames of a sequence were never optimized individually belongs to the problem of linear phase interpolation, whose frames are not optimized~\cite{lpi}, and not to the field. Knottnerus et al.\ describe earlier rearrangement as holograms calculated by the weighted algorithm with free phases that jump~\cite{lpi}; the induction of Kim et al.\ is the published answer to the jump, and its measured step of \accKimFiveStepMax{} rad shows that the answer works for the task of this paper. Linear phase interpolation of Knottnerus et al.\ is a different trade: it is cheap, it has no iterations in the frames, and its frames are not uniform, which is the reason it dips.

\subsection{What the results imply for use}

For a move that fits in a small number of frames and an experiment that tolerates the phase steps of the induced sequence, per-frame Gerchberg--Saxton with phase induction reads as better or tied on the dip in every cell of this paper and costs half the seconds at its own iteration count. The proposed sequence is the better choice where a bound on the phase step between frames is a requirement, for instance when the phase slip per hologram is the measured cause of loss, or where the per-frame uniformity matters more than the dip; there the lock provides the bound as a parameter and not as an empirical outcome. For a long move the frame count of Eq.~\ref{eq:frames} is a starting point, to be verified on the seeds of the experiment, and the number of iterations per frame can be lowered. A linear profile is the better setting for the dip with the same frames. Linear phase interpolation is the choice when computation time per frame is the constraint, since it has no iterations in the frames, and the loss in the dip is the price.
